\documentclass{article}
\usepackage{mathnotes}

\title{Series}
\description{Studies infinite series convergence through partial sums, covering convergence tests including geometric series, p-series, ratio test, root test, and comparison tests.}

\begin{document}

\section{Series}

An infinite sequence $a_n$ can be used to form an infinite series, or simply series, by summing its entries:

\[ \sum_{n=1}^{\infty} a_n = a_1 + a_2 + a_3 + \cdots \]

We can define a sequence $s_n$ of partial sums of the infinite series $\sum_{n=1}^{\infty}{a_n}$ by assigning

\[ s_n = \sum_{k=1}^{n} a_n = a_1 + a_2 + \cdots + a_n. \]

The series $\sum_{n=1}^{\infty}{a_n}$ is said to be convergent and to have a sum $m$ if the sequence of partial sums $s_n$ converges to $m.$ If the sequence $s_n$ diverges, then the series is said to diverge.

\subsection{Absolute Convergence}
The series $\sum_{n=1}^{\infty}{a_n}$ is said to \textbf{converge absolutely} if $\sum_{n=1}^{\infty}{\|a_n\|}$ converges. The infinite series $\sum_{n=1}^{\infty}{a_n}$ is said to \textbf{converge conditionally} if it converges but $\sum_{n=1}^{\infty}{\|a_n\|}$ diverges.

\subsection{Linearity of Summation of Convergent Series}

\emph{Theorem}: If $c$ is a \@{real-number} and the series $\sum_{n=1}^{\infty}{a_n}$ and $\sum_{n=1}^{\infty}{b_n}$ are convergent, then

\[ \sum_{n=1}^{\infty}{ \left ( a_n + c \cdot b_n \right )} =  \sum_{n=1}^{\infty}{a_n} + c \sum_{n=1}^{\infty}{b_n}. \]

\subsection{Convergence Tests}

Before considering the sum of a series, we need to know if it converges or not. There are many tests for convergence; we will start with a test for divergence.

\subsubsection{Divergence Test}

\emph{Theorem}: The \textbf{divergence test} says that if an infinite series $\sum_{n=1}^{\infty}{a_n}$ converges, then $\lim_{n \to \infty}{a_n} = 0.$ Therefore, if $\lim_{n \to \infty}{a_n} \neq 0$ (if the limit diverges or converges to a value other than 0), the series diverges.

\emph{Proof}: Suppose $\sum_{n=1}^{\infty}{a_n}$ converges with sum $m$ and let $s_n = \sum_{k=1}^{n} a_n.$ Then, it follows that

\[ \lim_{n \to \infty} a_n = \lim_{n \to \infty}{(s_n - s_{n-1})} = \lim_{n \to \infty}{s_n} - \lim_{n \to \infty}{s_{n-1}} = m - m = 0. \square \]

Note that $a_n$ converging to $0$ is a necessary condition for $\sum_{n=1}^{\infty}{a_n}$ to converge, but it is not sufficient.

\subsubsection{Geometric Series Test}

\emph{Theorem}: The \textbf{geometric series test} says that the geometric series $\sum_{n=1}^{\infty}{r^n}$ converges if $\|r\| < 1$ and diverges otherwise. If $\|r\| < 1,$ then $\sum_{n=1}^{\infty}{r^n} = \frac{r}{1 - r}.$

\emph{Proof}: Suppose $\|r\| < 1.$ Note that

\[ s_n = (r + r^2 + \cdots + r^n) = r(1 + r + r^2 + \cdots + r^{n-1}) = r\frac{1-r^n}{1 - r}. \]

(That last step is a rabbit out of a hat but can be shown using long division.) Now, if we take the limit as n goes to infinity we have

\[ \lim_{n \to \infty}{s_n} = \lim_{n \to \infty}{r \frac{1 - r^n}{1 - r}} = r \frac{1}{1 - r} = \frac{r}{1 - r}. \square \]

Note that when $\|r\| \geq 1,$ the series diverges.

\subsubsection{Bounded Series with Nonnegative Terms}

\emph{Theorem:} Suppose $a_n \geq 0$ for all $n \in \naturals.$ Then the series $\sum_{n=1}^{\infty}{a_n}$ converges iff the sequence of partial sums $s_n = \sum_{k=1}^{n}{a_k}$ is bounded.

\emph{Proof:} Assume the sequence of partial sums $s_n$ is bounded. Since all terms of $a_n$ are nonnegative, $s_n$ is increasing, and by the Monotone Convergence Theorem, must converge, meaning $\sum_{n=1}^{\infty}{a_n} = \lim_{n \to \infty} s_n$ converges. On the other hand, assume $\sum_{n=1}^{\infty}{a_n}$ converges. Then $s_n$ is a convergent sequence and is therefore bounded (see \pagelink[sequences]{analysis/real-analysis/real-sequences}).

\subsubsection{$2^n$ Test}
\emph{Theorem:} Suppose $a_n \geq a_{n+1} \geq {0}$ for all $n \in \naturals$ (i.e. $a_n$ is nonnegative and decreasing.) Then the series $\sum_{n=1}^{\infty}{a_n}$ converges iff the series $\sum_{n=1}^{\infty}{2^n a_{2n}}$ converges.

\subsubsection{$p$-Series Test}
\emph{Theorem:} The $p$-series $\sum_{n=1}^{\infty}{\frac{1}{n^p}}$ converges if $p > 1$ and diverges otherwise.

\subsubsection{Alternating Series Test}
\emph{Theorem:} Suppose $a_n \geq a_{n+1}$ for all \@{natural-numbers} $n$ and that $\lim_{n \to \infty}{a_n} = 0.$ Then the alternating series $\sum_{n=1}^{\infty}{(-1)^{n+1}a_n}$ converges.

\subsubsection{Absolute Value Test}
\emph{Theorem:}  If the series $\sum_{n=1}^{\infty}{a_n}$ converges absolutely then it simply converges too.

\subsubsection{Comparison Test}
\emph{Theorem:} Suppose $\sum_{n=1}^{\infty}{a_n}$ and $\sum_{n=1}^{\infty}{b_n}$ are series for which $\|a_n\| \leq \|b_n\|$ for all natural $n.$ Then if $\sum_{n=1}^{\infty}{b_n}$ converges absolutely, so does $\sum_{n=1}^{\infty}{a_n}.$ If $\sum_{n=1}^{\infty}{\|a_n\|}$ diverges, so does $\sum_{n=1}^{\infty}{\|b_n\|}.$

\subsubsection{Limit Ratio Test}
\emph{Theorem:} For the series $\sum_{n=1}^{\infty}{a_n},$ let

\[ L = \lim_{n \to \infty}{\left | \frac{a_{n+1}}{a_n} \right |}. \] 

Then, if $L < 1,$ $\sum_{n=1}^{\infty}{a_n}$ converges absolutely. If $L > 1,$ $\sum_{n=1}^{\infty}{a_n}$ diverges.

\subsubsection{Root Test}
\emph{Theorem:}  For the series $\sum_{n=1}^{\infty}{a_n},$ let $L = \limsup_{n \to \infty}{\|a_n\|^{1/n}}.$ Then if $L < 1,$ $\sum_{n=1}^{\infty}{a_n}$ converges absolutely. If $L > 1,$ $\sum_{n=1}^{\infty}{a_n}$ diverges.

\end{document}
